% =====================================================================
%  CHAPTER 8: 儿童少年常见病（对应大纲第七章）
%  参考PPT：第八章 儿童少年常见疾病（邹志勇，北京大学）
%  往届笔记：第九章 儿童少年常见病
% =====================================================================
\chapter{儿童少年常见病}
\setcounter{chapter}{8}
\chapterintro{
\textbf{【熟悉/单横线】}健康的含义：身体、心理、社会适应和道德上的完好状态。

\textbf{【掌握/双横线】}儿童少年常见病的概念和流行特点；单纯性肥胖的定义、发病特点、流行规律及防控措施；视力低下（近视）的定义、分类、发生机制、流行规律及防控措施；龋齿和牙周病的流行状况、四联致病因素及预防控制；脊柱弯曲异常的流行特点、发生原因及预防控制。
}

% ===== 8.1 概述 =====
\section{概述}

\subsection{儿童少年常见病概念与流行特点}

\begin{defbox}
\textbf{儿童少年常见病（Common Diseases in Children and Adolescents）：}在儿童少年中发病率较高的一类疾病，发病与其正处于生长发育阶段有关，也受幼儿园和学校集体生活影响。
\end{defbox}

\begin{keybox}
\textbf{中国儿童少年常见病流行特点：}
\begin{enumerate}[leftmargin=*]
    \item \imp{发病率基本得到控制的常见病}：沙眼、蛔虫感染
    \item \imp{患病率居高不下的常见病}：近视
    \item \imp{检出率持续上升的常见病}：超重肥胖
    \item \imp{发展不平衡造成的常见病}：营养不良
\end{enumerate}
\end{keybox}

\subsection{危险因素}

\begin{keybox}
\begin{enumerate}[leftmargin=*]
    \item \imp{遗传因素}
    \begin{itemize}[leftmargin=*]
        \item 单基因遗传：高度近视，常染色体显性、常染色体隐性、性连锁隐性遗传，有家族聚集性
        \item 多基因遗传：超重肥胖、单纯性近视；双生子研究显示，近视遗传度为65\%$\sim$70\%
    \end{itemize}
    \item \imp{环境因素}：环境内分泌干扰物（EDCs），家庭环境（生活环境、家庭结构、家庭经济状况、父母文化程度、教育方式）
    \item \imp{个体因素}：膳食摄入、饮食习惯、体力活动
\end{enumerate}
\end{keybox}

\subsection{预防控制}

\begin{keybox}
\begin{enumerate}[leftmargin=*]
    \item \imp{健康促进}：(1)制定并认真落实学校健康促进政策；(2)保证学生每天体育锻炼、户外活动时间；(3)定期组织体检，建立学生健康档案
    \item \imp{环境改善}：(1)改善教室环境；(2)创造良好家庭和学习环境
    \item \imp{健康素养提升}：(1)提高对常见病认知程度和防范意识；(2)养成良好行为和生活方式
\end{enumerate}
\end{keybox}

% ===== 8.2 肥胖 =====
\section{儿童少年肥胖}

\begin{defbox}
\textbf{肥胖（Obesity）：}在遗传和环境因素交互作用下，因能量摄入$>$能量消耗，导致体内脂肪积聚过多，危害健康的慢性代谢性疾病。

\begin{itemize}[leftmargin=*]
    \item WHO标准：超重25$\leq$BMI$\leq$30，肥胖BMI$\geq$30
    \item 中国标准：超重24$\leq$BMI$\leq$28，肥胖BMI$\geq$28
\end{itemize}
\end{defbox}

\subsection{流行特点}

\begin{keybox}
\textbf{全球：}1975年儿童青少年肥胖率$<$1\%，2022年上升至8\%（女性6.9\%，男性9.3\%）。

\textbf{中国：}7$\sim$18岁儿童青少年肥胖率1985年0.1\% $\rightarrow$ 2019年9.6\%（增长75.6倍），预计2030年将达到15.1\%。农村增长速度加快，呈现"城乡逆转"趋势。
\end{keybox}

\subsection{危险因素}

\begin{keybox}
\textbf{继发性肥胖：}病因明确，与原发性病因有关，需先治疗其原发性病因。

\textbf{单纯性肥胖（敏感期：孕后期、婴儿期、青春早期、青春后期，无幼儿期和青春中期）}

\begin{enumerate}[leftmargin=*]
    \item \imp{遗传因素}
    \begin{itemize}[leftmargin=*]
        \item 是重要因素，但不是决定因素，肥胖是多基因关联累加结果，属于多基因遗传
        \item 瘦素通过3种途径调节机体脂肪沉积：①抑制食欲；②增加能量消耗；③抑制脂肪合成
    \end{itemize}
    \item \imp{个人因素}
    \begin{itemize}[leftmargin=*]
        \item 膳食摄入：A.膳食热能过高；B.膳食结构不合理；C.食用能量密度高的食物
        \item 饮食行为：A.不良食品选择倾向——甜食、甜饮料等；B.不良饮食方式——煎、炸烹饪方式等；C.不良膳食制度——不吃早餐等；D.不良食物奖惩形式——吃西式快餐作为奖励
        \item 体力活动：A.体育活动$\downarrow$；B.静态生活时间$\uparrow$
        \item 肠道菌群
    \end{itemize}
    \item \imp{环境因素}
    \begin{itemize}[leftmargin=*]
        \item 社会环境因素：A.传统健康观念——"胖是健康，以胖为福"；B.肥胖易感环境（Obesogenic Environment）——全球化、城市化带来的膳食快餐化、强大市场诱惑、媒体对饮食行为的影响
        \item 家庭环境因素：受母亲文化程度等因素影响，文化程度高的母亲更倾向于调整孩子的饮食量，鼓励子女参加体育活动
    \end{itemize}
\end{enumerate}
\end{keybox}

\subsection{肥胖危害}

\begin{keybox}
\begin{enumerate}[leftmargin=*]
    \item \imp{健康危害}：心血管疾病、早期动脉粥样硬化、2型糖尿病、代谢综合征、非酒精性脂肪肝、阻塞性睡眠呼吸暂停（OSA）
    \item \imp{心理危害和行为危害}
    \begin{itemize}[leftmargin=*]
        \item 心理危害：心理压抑，心理精神障碍（如抑郁和自杀意念等）
        \item 行为危害：饮食失调症、社会交往能力$\downarrow$、吸烟饮酒行为、校园欺负行为，甚至自杀意念和行为
    \end{itemize}
    \item \imp{社会危害}
    \begin{itemize}[leftmargin=*]
        \item 直接成本：控制和治疗肥胖付出的成本
        \item 机会性成本：肥胖相关疾病或过早死亡造成的医疗成本或经济损失
        \item 间接成本：个人和社区的间接（社会）负担，如病假、个人用于减轻体重的花费等
    \end{itemize}
\end{enumerate}
\end{keybox}

\subsection{肥胖筛查}

\begin{defbox}
\textbf{筛查方法：}
\begin{enumerate}[leftmargin=*]
    \item 目测法
    \item 身高标准体重法：WHO推荐使用，限于0$\sim$6岁儿童
    \item 体质量指数（BMI）：简便易测，能进行跨人群（现状）和跨时段（趋势）分析
    \item 腹部脂肪测量：对预测肥胖造成的危害和儿科临床应用有重要意义和价值
\end{enumerate}

\textbf{儿童少年肥胖筛查标准：}
\begin{table}[h]
\centering
\begin{tabular}{lccc}
\hline
标准 & NCHS标准 & IOTF标准 & 中国标准 \\
\hline
超重 & P85$\leq$BMI$\leq$P95 & 25$\leq$BMI$\leq$30 & 24$\leq$BMI$\leq$28 \\
肥胖 & BMI$\geq$P95 & BMI$\geq$30 & BMI$\geq$28 \\
\hline
\end{tabular}
\end{table}

\textbf{体成分测定：}体脂率（Body Fat Percentage, \%BF）：人体脂肪组织占体重百分比，是较直观判断肥胖的指标。
\end{defbox}

\subsection{预防控制}

\begin{keybox}
\begin{enumerate}[leftmargin=*]
    \item \imp{普遍性预防}：运用健康教育、健康促进理论，依托建设健康学校、提高学生健康素养的形式，通过制定政策、创建支持环境、社区积极参与、开展健康教育、提供健康咨询和指导，培养儿童少年健康行为习惯和生活方式
    \item \imp{针对性预防}
    \begin{itemize}[leftmargin=*]
        \item 平衡膳食、建立良好膳食制度、保持食物多样化、三餐热量合理分配
        \item 培养健康饮食行为，少吃油炸食品、不喝含糖饮料等
        \item 积极参加体育活动（每天60min$\uparrow$中高等强度体育活动），减少静态活动时间（每天看电视玩游戏$\leq$2h）
        \item 不提倡节食、挑食，不使用减肥药、腹泻等不健康"减肥"方式
    \end{itemize}
    \item \imp{干预性防控}
    \begin{itemize}[leftmargin=*]
        \item 饮食调整：膳食结构合理，减少脂肪摄入
        \item 身体活动指导：坚持规律有氧运动，通过评估调整运动方案
        \item 行为矫正：制定行为改变目标，以日记等方式进行自我监督，正向鼓励良好行为，家长以身示范
        \item 心理疏导：树立正常健康观和意识，疏解压抑、自卑、消极心理
    \end{itemize}
\end{enumerate}
\end{keybox}

% ===== 8.3 近视 =====
\section{儿童少年近视}
\important{中小学生近视筛查频率：每年2次。}

\begin{defbox}
\textbf{近视：}不使用调节功能状态下，远处来的平行光线在视网膜感光层前方聚焦。

\textbf{诊断标准：}
\begin{itemize}[leftmargin=*]
    \item 近视：等效球镜度数 $\leq -0.50$D
    \item 高度近视：等效球镜度数 $\leq -6.0$D（600度及以上）
\end{itemize}

\textbf{分类：}
\begin{enumerate}[leftmargin=*]
    \item 按屈光度：低度近视（$-0.25\sim-3.0$D），中度近视（$-3.25\sim-6.0$D），高度近视（$-6.25\sim-9.0$D）
    \item 按有无调节因素参与：假性，真性，半真性
    \item 按屈光要素改变：轴性，屈光性
\end{enumerate}
\end{defbox}

\subsection{流行特点}

\begin{keybox}
\begin{itemize}[leftmargin=*]
    \item 中国是全球近视患病率最高的国家之一
    \item 随年龄/年级增长患病率持续上升（年级效应$>$年龄效应）
    \item 城市高于农村，但农村增长速度更快
    \item 女性高于男性
    \item 2010年、2014年、2019年全国学生体质调研数据显示，近视检出率持续走高，低龄化趋势明显
\end{itemize}
\end{keybox}

\subsection{发生机制}

\begin{keybox}
\textbf{正视化（Emmetropization）：}大部分儿童出生时远视，随发育进程眼轴变长、晶状体和角膜弯曲度变平、眼睛屈光度变强，发展为正视的过程。

\textbf{生理基础：}正视化完成后角膜曲率基本不再变化，但眼轴长度可继续生长。

学习时间过长，光照条件不良$\rightarrow$眼睛常处于高度调节紧张状态$\rightarrow$看远处时睫状肌仍处于收缩状态$\rightarrow$悬韧带放松$\rightarrow$晶状体屈光力过强$\rightarrow$近视。

\begin{itemize}[leftmargin=*]
    \item \imp{假性近视（Pseudomyopia）}：因过度视近工作形成的近视为调节紧张性近视，属功能性改变，是近视防控关键窗口期。采取视力保护措施$\rightarrow$纠正不良用眼习惯，放松睫状肌$\rightarrow$视力可恢复正常
    \item \imp{真性/轴性近视}：假性近视后依然不注意用眼卫生$\rightarrow$引起眼球充血、眼压升高、眼球壁弹性下降$\rightarrow$刺激眼轴伸长
\end{itemize}
\end{keybox}

\subsection{影响因素}

\begin{keybox}
\begin{enumerate}[leftmargin=*]
    \item \imp{遗传因素}
    \item \imp{环境因素}
    \item \imp{体质健康因素}
\end{enumerate}
\end{keybox}

\subsection{危害}

\begin{keybox}
\begin{itemize}[leftmargin=*]
    \item 远视力下降（最典型表现）
    \item 视疲劳
    \item 夜间视力差、飞蚊症（高度近视）
    \item 视网膜脱离/裂孔、黄斑出血、青光眼风险增加（高度近视）
    \item 严重者可致盲
\end{itemize}
\end{keybox}

\subsection{预防控制}

\begin{keybox}
\begin{enumerate}[leftmargin=*]
    \item \imp{近视预防}：(1)加强学校预防近视工作；(2)培养良好用眼习惯；(3)创造良好生活环境；(4)认真做好眼保健操；(5)积极参加户外活动，亲近阳光
    \item \imp{近视矫正}：(1)及时检查；(2)合理配镜；(3)抗胆碱能药物；(4)手术治疗
\end{enumerate}
\end{keybox}

% ===== 8.4 龋齿和牙周病 =====
\section{龋齿和牙周病}
\important{中小学生龋齿筛查频率：每年1次。}

\begin{defbox}
\textbf{龋齿（Dental Caries）：}牙齿在身体内外因素作用下，硬组织脱矿，有机质溶解，牙组织进行性破坏，导致牙齿缺损的儿童常见病。
\end{defbox}

\subsection{流行状况与指标}

\begin{keybox}
\textbf{【考点】$\triangle$流行状况：}龋患率；\textbf{【考点】$\triangle$严重程度：}龋均、患者龋均；\textbf{【考点】$\triangle$治疗状况：}DMFT。
\end{keybox}

\subsection{龋齿发生原因与危害}

\begin{keybox}
\textbf{发生原因——四联致病因素：}
\begin{enumerate}[leftmargin=*]
    \item \imp{细菌和菌斑}
    \item \imp{宿主因素}
    \item \imp{食物因素}
    \item \imp{时间因素}
\end{enumerate}

\textbf{危害：}
\begin{enumerate}[leftmargin=*]
    \item \imp{局部}：影响食欲，干扰咀嚼、消化和吸收，导致营养缺乏，且随龋病发展可引发牙髓炎、颜面蜂窝织炎、根周脓肿、齿槽溢脓等严重口腔疾病
    \item \imp{全身性疾病}：通过变态反应引发风湿性关节炎、心内膜炎、肾炎
    \item \imp{损害个人形象}，影响心理健康
\end{enumerate}
\end{keybox}

\subsection{牙周病}

\begin{defbox}
\textbf{分类：}
\begin{enumerate}[leftmargin=*]
    \item 牙龈病：病变仅累及牙龈组织
    \item 牙周炎：病变累及牙龈组织$+$深层组织
\end{enumerate}

\textbf{发生原因：}细菌，宿主因素，环境因素。
\end{defbox}

\subsection{预防控制}

\begin{keybox}
\begin{enumerate}[leftmargin=*]
    \item \imp{加强口腔保健}：定期检查，早期诊断
    \item \imp{控制牙菌斑}：刷牙3min
    \item \imp{注意饮食卫生，增强宿主抗龋力}
    \item \imp{加强健康教育}
    \item \imp{健全学校口腔疾病防治网}
\end{enumerate}
\end{keybox}

% ===== 8.5 脊柱弯曲异常 =====
\section{脊柱弯曲异常}

\begin{defbox}
\textbf{脊柱弯曲异常（Vertebral Column Defects）：}脊柱弯曲超出正常生理范围。
\end{defbox}

\subsection{流行特点}

\begin{keybox}
2019年起脊柱弯曲异常纳入全国学生常见病监测。全国初筛阳性率约为2\%$\sim$3\%。
\end{keybox}

\subsection{发生原因}

\begin{keybox}
\begin{enumerate}[leftmargin=*]
    \item \imp{习惯性姿势}：不良站姿，不良坐姿，歪头写字
    \item \imp{桌椅高矮不适合}
    \item \imp{体力活动缺乏}
    \item \imp{营养和体质因素}
\end{enumerate}
\end{keybox}

\subsection{预防控制}

\begin{keybox}
\begin{enumerate}[leftmargin=*]
    \item 加强健康教育
    \item 保证足够体育锻炼和体力活动
    \item 调整座椅高度
    \item 定期筛查，早期发现
    \item 脊柱弯曲异常矫正
\end{enumerate}
\end{keybox}

% ===== 营养不足 =====
\section{营养不足}

\subsection{流行特点}

\begin{defbox}
\textbf{营养不足（undernutrition）的三种表现：}

\begin{enumerate}[leftmargin=*]
    \item \imp{生长迟缓}：表现为年龄别身高低，发育迟缓会阻碍儿童发挥其身体和认知潜力。
    \item \imp{消瘦}：表现为身高别体重低，中度或重度消瘦的年幼儿童面临更高死亡风险。
    \item \imp{体重不足}：表现为年龄别体重低，体重不足的儿童可能会生长迟缓、消瘦或两者兼有。
\end{enumerate}
\end{defbox}

\begin{keybox}
\textbf{流行现状：}
中国曾是全球儿童营养不足最高发的国家之一，新中国成立初约50\%的儿童营养不足。中国中小学生营养不足率从2010年的12.7\%下降到2019年的8.5\%，下降了4.2个百分点。《中国儿童发展纲要（2011—2020年）》提出要控制中小学生营养不足发生率。
\end{keybox}

\subsection{发生原因}

\begin{keybox}
\begin{enumerate}[leftmargin=*]
    \item \imp{营养素摄入不足}：儿童少年对热能、优质蛋白质的需求若得不到足量供应，将导致营养不足发生。
    \item \imp{疾病}：肠道寄生虫感染、肺结核、肝炎、胃病等主要脏器慢性疾病，是造成营养不足的最常见疾病。
    \item \imp{膳食结构和饮食习惯}：不良膳食食物种类单一，结构不合理；一日三餐热量、营养素搭配不合理；饮食习惯不良。
    \item \imp{早期生长潜力未充分发挥}：婴幼儿阶段是决定成年身高的关键时期，儿童身高增长的最大潜能时期是在6个月至3岁期间。
    \item \imp{体像认知和心理因素}：片面追求"体型美"，盲目节食。
\end{enumerate}
\end{keybox}

\subsection{危害}

\begin{keybox}
\begin{itemize}[leftmargin=*]
    \item \imp{生长迟缓}：阻碍儿童发挥其身体和认知潜力
    \item \imp{消瘦}：中度或重度消瘦增加死亡风险
    \item \imp{体重不足}：可兼有生长迟缓和消瘦两者的不良后果
\end{itemize}
\end{keybox}

\subsection{预防控制}

\begin{keybox}
2014年联合国提出消除全球营养不良共识。2017年中国国务院发布《国民营养计划（2017—2030年）》，旨在降低5岁以下儿童贫血率并控制肥胖率上升趋势。针对中国儿童少年群体：
\begin{enumerate}[leftmargin=*]
    \item 采取家庭、社区和学校相结合的营养知识、行为习惯、生活方式宣教。
    \item 帮助青少年建立正确的体像观。
    \item 培养良好的饮食习惯和生活方式。
    \item 指导学生合理饮食，开展超重和肥胖干预。
\end{enumerate}
\end{keybox}

% ===== 缺铁性贫血 =====
\section{缺铁性贫血}

\subsection{流行特点}

\begin{keybox}
\textbf{全球：}40\%的6至59月龄儿童、37\%的孕妇和30\%的15～49岁妇女受到贫血影响。2019年，贫血导致因残疾损失了5000万健康生命年。最大原因是膳食中缺铁、地中海贫血和镰状细胞性状以及疟疾。

\textbf{中国：}2010年中国中小学生贫血率为11.3\%，2014年下降到9.3\%，但2019年又反弹至11.1\%。低年龄组小学生和青春期少年是两大易感人群。
\end{keybox}

\subsection{发生原因}

\begin{keybox}
\begin{enumerate}[leftmargin=*]
    \item \imp{体内需求量大}：儿童少年生长发育旺盛，尤其青春期生长突增以及组织代谢的氧需求和女童月经周期等铁的需求量很大。
    \item \imp{铁的额外丢失}：
    \begin{itemize}
        \item 月经不调性失血：是儿童少年铁额外丢失的首位原因。
        \item 外伤：引起急慢性失血的常见原因。
        \item 疾病：胃十二指肠溃疡、胃酸缺乏、服用抗酸性药物、肠道蠕虫感染、慢性腹泻。
        \item 环境污染：铅、苯、镉等中毒可大量破坏体内红细胞造成慢性失血。
    \end{itemize}
\end{enumerate}
\end{keybox}

\subsection{预防控制}

\begin{keybox}
\begin{enumerate}[leftmargin=*]
    \item \imp{去除病因}：食不当者应纠正不合理的饮食习惯和食物组成，有偏食习惯者应予纠正。如有器质性疾病，应积极治疗。
    \item \imp{一般治疗}：轻度贫血儿童少年应合理膳食、注意营养，增加富铁性食物的摄入；尽量纠正偏食、挑食及过度节食等不良饮食习惯。
    \item \imp{铁剂治疗}：在医生监督下合理补充铁剂可使血红蛋白尽快恢复正常，铁储备得到补充。
    \item \imp{防治结合的综合措施}：儿童青少年缺铁性贫血发病率高、治疗容易但复发率高，同时容易预防，生活中应采取综合措施进行防控。
\end{enumerate}
\end{keybox}

% ===== 哮喘 =====
\section{哮喘}

\subsection{流行特点}

\begin{keybox}
\textbf{全球：}2019年儿童哮喘死亡1.29万人，发病率2200万，DALY 510万（与1990年相比，上述指标均有所下降）。

\textbf{中国：}三次全国调查（1990年、2000年、2010年）显示0$\sim$14岁城市儿童累计患病率持续上升：1.09\% $\rightarrow$ 1.97\% $\rightarrow$ 3.02\%。
\end{keybox}

\subsection{病因与危险因素}

\begin{keybox}
\begin{enumerate}[leftmargin=*]
    \item \imp{遗传因素}：易感者的基因表达、免疫状态、精神状态、内分泌状态及健康状态等多因素协同作用，决定个体是否发生哮喘。
    \item \imp{环境因素}
    \begin{itemize}[leftmargin=*]
        \item 变应原（关键因素）：包括吸入性变应原（如尘螨、花粉、霉菌等）和食物性变应原（如牛奶、鸡蛋、海鲜等）
        \item 呼吸道病毒感染：是诱导气道高反应性的重要原因
        \item 气候变化
        \item 空气污染
        \item 食物过敏
        \item 其他：金首饰、香水等特殊物品
    \end{itemize}
\end{enumerate}

\textbf{分类：}
\begin{itemize}[leftmargin=*]
    \item \imp{外源性哮喘}：特应性/季节性，有明确过敏原触发
    \item \imp{内源性哮喘}：非特应性，无明确过敏原触发
\end{itemize}
\end{keybox}

\subsection{危害}

\begin{keybox}
\begin{itemize}[leftmargin=*]
    \item 突然反复发作的喘息、气短、胸闷、咳嗽
    \item 可变性呼气气流受限
    \item 运动后或夜间/凌晨发作加重
    \item 影响学习、生活质量和心理发展
\end{itemize}
\end{keybox}

\subsection{预防控制}

\begin{keybox}
\textbf{三级预防策略：}
\begin{enumerate}[leftmargin=*]
    \item \imp{一级预防}：避免接触致敏原，消除高危因素，从母孕期入手
    \item \imp{二级预防}：出现初发症状后加强预防，防止再次发作
    \item \imp{三级预防}：防止不可逆改变，避免后遗症
\end{enumerate}

\textbf{关键措施：}
\begin{itemize}[leftmargin=*]
    \item 避免接触变应原（关键）
    \item 普及防控知识，提高公众认知
    \item 适度锻炼，增强体质
\end{itemize}
\end{keybox}

\newpage
